% Source: https://artofproblemsolving.com/community/c3h1349077_sum_of_consecutive_numbers
% Stable identity: aops-topic-1349077-q01
A sum of $k\ge2$ consecutive positive integers starting at $a\ge1$ is $ka+\dfrac{k(k-1)}{2}$, so we need
\[
k(2a+k-1)=210 .
\]
Here $k<2a+k-1$ (since $a\ge1$), and exactly one of the two factors is even, because $210=2\cdot3\cdot5\cdot7$ contains a single factor $2$ while the two factors differ by the odd number $2a-1$. Conversely, every factorisation $210=k\cdot L$ with $2\le k<L$ gives $a=\dfrac{L-k+1}{2}$, a positive integer because $L-k$ is odd. The $16$ divisors of $210$ give $8$ factorisations with $k<L$; removing $k=1$ leaves $7$:
\[
\begin{aligned}
k=2&:\ 52+53, & k=3&:\ 34+35+36, & k=5&:\ 19+\dots+23, & k=6&:\ 15+\dots+20,\\
k=7&:\ 12+\dots+18, & k=10&:\ 6+\dots+15, & k=14&:\ 1+\dots+14 .
\end{aligned}
\]
So $105$ can be written in exactly $\boxed{7}$ ways, and eight ways are impossible. (Counting the one-term ``sum'' $105$ itself as well would give $8$.)
